\hypersetup{pageanchor=false}
\pagenumbering{gobble}
\begin{titlepage}
\begin{center}
Национальный исследовательский университет\\
«Высшая школа экономики»

\vfill
{\Large\textbf{Адаптационный курс}}\par
\vspace{0.7cm}
{\large\textbf{Практикум для подготовки\\к первой контрольной работе}}\par
\vspace{0.7cm}
Лекции и семинары №1--3\par
Системы счисления и представление чисел

\vfill
Для студентов первого курса направления\\
«Программная инженерия»

\vspace{1cm}
Москва, 2026
\end{center}
\end{titlepage}
\clearpage
\hypersetup{pageanchor=true}
\pagenumbering{arabic}

\chapter{Материалы для подготовки}
Примеры П1--П10 и задания З1--З32 относятся к занятиям №1--2. Контрольные К1--К8 и В1--В8 охватывают занятия №1--2 и №1--3 соответственно. Каждый вариант рассчитан на 90 минут и сопровождается ответами и подробными решениями.

Теория и задачи семинаров: \href{adaptation_course.pdf}{«Адаптационный курс»}, глава~1.
\begin{itemize}
    \item \href{adaptation_course.pdf}{\textbf{Лекция и семинар №1.}} Системы счисления, переводы и арифметика.
    \item \href{adaptation_course.pdf}{\textbf{Лекция и семинар №2.}} Коды и диапазоны целых чисел, побитовые операции, сдвиги и флаги.
    \item \href{adaptation_course.pdf}{\textbf{Лекция и семинар №3.}} Числа с плавающей точкой: IEEE 754, арифметика, округление и машинная точность.
\end{itemize}

\section{Обозначения и общие правила}
Основание указано индексом; \texttt{0xA9} означает $A9_{16}$.

\begin{itemize}
    \item Если ширина слова не указана, разрядность не ограничена.
    \item В целочисленных операциях сохраняются $n$ младших битов. Ведущие нули записываются до заданной ширины слова.
    \item Для целых чисел «знаковое» означает дополнительный код, если не указан другой формат. Для кода со смещением величина bias задана явно.
    \item Биты нумеруются справа с нуля.
    \item В таблицах флаги идут в порядке \textbf{CF, OF, ZF, SF}. Для вычитания используются правила инструкции \texttt{SUB} архитектуры x86.
\end{itemize}

\tableofcontents
